% !TEX root = ../supp.tex
\section*{B. Next Action Anticipation}
We conduct an experiment of next action anticipation on EK55~(validation, RGB) following the previous experimental protocols~\cite{girdhar2021anticipative,sener2020temporal,furnari2019would}.
\\
\noindent
\textbf{Dataset.}
The Epic-Kitchens 55 dataset~\cite{Damen2018EPICKITCHENS} is the large-scale dataset in first-person vision. The dataset comprises of 55 hours of recordings of 32 kitchens, including 39,594 action segments annotated with 125 verb, 331 noun, and 2,513 action classes. 
\\
\noindent
\textbf{Implementation details.}
\ours can be applied to next action anticipation by simply setting the number of action query $M$ to 1.
We use two encoder layers and two decoder layers while setting the size of the hidden dimension $D$ to 512.
We do not include action segmentation loss in this experiment due to the
lack of frame-level action annotations. 
Instead, we use additional a fully-connected layer applying to the output of the encoder layers $X_{L^E}$ to predict features of the next frame. Then we apply a feature prediction loss of L2 distance between predicted features and the next frame similar to AVT~\cite{girdhar2021anticipative}.
We use AdamW optimizer~\cite{loshchilov2017decoupled} with a learning rate of 1e-5. We train our model for 40 epochs with a batch size of 32.
We use the RGB feature embedded by TSN~\cite{wang2016_TemporalSegmentNetworks} in this experiment. 
\\
\noindent
\textbf{Results.}
\begin{table}[t]
\begin{center}
\scalebox{1.0}{
\begin{tabular}{lcc}
        \toprule
        method & backbone & top-1 \\
        \cmidrule{1-3}
        RULSTM~\cite{furnari2019would}      & TSN & 13.1  \\
        Temporal Agg.~\cite{sener2020temporal} & TSN & 12.3  \\
        AVT~\cite{girdhar2021anticipative}    & TSN & 13.1 \\
        FUTR (ours) & TSN & 12.3  \\
        \bottomrule 
        \end{tabular}
        }
\vspace{-1mm}
\counterwithin{table}{section}
\renewcommand\thetable{S\arabic{table}}
\caption{\textbf{Performance comparison on EK55.} Although FUTR is designed for long-term action anticipation, the model is also effective in next action anticipation.}
\label{tab:EK55}
\end{center}
\vspace{-4mm}
\end{table}
% \setlength\intextsep{1pt}
% \begin{wraptable}{l}{0pt}
% \centering
% \setlength{\tabcolsep}{2pt}
% \resizebox{0.17\textwidth}{!}{
% \begin{tabular}{ccc}
%         \cmidrule[0.8pt]{1-3}
%         model & input & top-1 \\
%         % \cmidrule[0.5pt]{1-4}
%         % Time-Cond.[19] & I3D & 25.6 & 71.6 \\
%         \cmidrule[0.5pt]{1-3}
%         RULSTM [14]      & TSN & 13.1  \\
%         Temporal Agg.[31] & TSN & 12.3  \\
%         AVT [17]    & TSN & 13.1 \\
%         FUTR (ours) & TSN & 12.3\\
%         % FUTR (ours) & TSN & 12.x  \\
%         \cmidrule[0.8pt]{1-3} 
%         \end{tabular}
% \hspace*{-1.2\columnsep}
% }
% \vspace{-3mm}
% \centering
% \caption{\textbf{Top-1 on EK55}}
% \label{tab:EK55}
% \end{wraptable}
% wl{3cm}C{2cm}C{2cm}
The result is shown in Table~\ref{tab:EK55}. \ours obtains 12.3\%p at top-1 accuracy performing comparable with the state-of-the-art methods. We find that \ours is also effective for next action anticipation, although the model is designed for long-term action anticipation. 

